\phantomsection
\addcontentsline{toc}{section}{Abstract}
\begin{center}\bfseries Abstract\end{center}
\begingroup
\fontsize{9.5}{10.7}\selectfont
\setlength{\parskip}{1.5pt plus .5pt minus .5pt}
\noindent
We derive the Bose--Einstein and Fermi--Dirac distributions, their dilute
Maxwell--Boltzmann limit, and the radiation laws of Planck,
Stefan--Boltzmann, and Wien from the quantum and thermal realizations
constructed within Triadic Modular Holography and Dimensional Mechanics.
The arithmetic path support, termed \emph{All Possible Paths} (APP),
is the discrete torus \((\mathbb Z/9\mathbb Z)^2\), equipped with an
additive Cayley graph in the cardinal directions and two evaluations,
addition and multiplication. The oriented ternary signature, TRIT, takes
values in \(\{-1,0,+1\}\) and distinguishes orientation and regime;
the \emph{Topological Phase Kernel} (TPK) composes selection, transport,
update, and prolongation while retaining memory. This dynamics produces
the regional states and records on which unitary evolution and the
calendar action are defined.

Parity determines the symmetric and exterior completions of these
realizations: the former admits nonnegative integer occupations; the
latter yields algebraic Pauli exclusion. The multiplicities determine the microcanonical counts;
Gibbs-state factorization gives the two equilibrium distributions,
whose dilute limit is quantified by explicit bounds. For free
transverse electromagnetic modes with zero chemical potential, bosonic
occupation and mode density produce the Planck spectrum. Convergence of
the discrete sums is proved through uniform bounds at low and high
frequencies. Their moments determine number, energy, and entropy
densities; angular integration gives the Stefan--Boltzmann flux, and
extremization of the frequency and wavelength densities gives the two
Wien laws.

The calendar and holonomy also produce a capacity renewal
$(20,24,25,23)$, a graded carrier with multiplicities $1/380/649$, and a
native geometric spectrum of ratio $1/10$. Its tritic refinement retains the
residue and gives $1/1000$; the marked $23/26/29$ section evaluates exactly
$1.380649\times10^{-23}$. Explicit isometries transport the energy origin,
marker, and provenance. Boundary currents fix the sector offset, while
composition with action proves
$G_a k_B\Theta_{{\rm clk},a}\log3=c^4L_\alpha$ and the effective transducer
identity on its positive section.

Configuration entropy, the measure on histories, and the spectrum of a
reduced state are related through transport and reduction maps. A
relative-divergence identity characterizes the stationary measure of
maximal entropy on the areal--volumetric envelope. Thermal transduction
connects information, action, and calendar with their dimensional bases.
The representation, equilibrium, and convergence conditions specify the
domain of each result.

In the finite realization \(\Lambda=(\mathbb Z/27\mathbb Z)^2\), the
unitary Fourier transform establishes
\(|\operatorname{supp}\psi|\,|\operatorname{supp}\widehat\psi|\ge729\)
and \(H_{\rm pos}(\psi)+H_F(\psi)\ge\log729\); four subgroup collections
saturate both bounds. On the dodecaphasic domain, the F044 readers
distinguish ordered correlations not retained by marginal entropy. The
prime-vacancy observer \(R_{\rm vac}\) composes Sturmian incidence with
the logarithmic clock and satisfies a modal identity under symmetric or
Fejér summation. Its six cutoffs through \(10^8\) and the thirty recorded
modes are finite witnesses; independent reproduction reaches \(10^6\).
The critical asymptotic estimate requires modal bounds uniform in \(k\).

\medskip
\noindent\textbf{Keywords:} Bose--Einstein; Fermi--Dirac; Planck law;
genealogical states; Fock completion; Floquet dynamics; occupation entropy;
modular Hamiltonian; thermodynamic limit; thermal radiation; entropic
uncertainty; prime vacancies; ordered memory.
\par\endgroup
